\section{Exclusion, Entry, and Price Formation}\label{sec:price_formation}

The set-aside changes two margins at once. It removes non-SMEs from eligible
auctions, but it also attracts more SMEs into the protected pool. This section
uses the counterfactuals from Section~\ref{sec:model} to separate the two
forces. The result is simple: SME entry responds strongly, but the price effect
is dominated by the loss of non-SME competitive discipline.

\subsection{Entry responds}

The first-stage participation facts already rule out a simple story in which
prices rise because protected firms fail to enter. In non-pharmaceutical
Preg\~ao auctions, the average number of SME bidders rises from
\bneNsmePreNp{} to \bneNsmePostNp{} per auction after the cutoff. In
pharmaceutical auctions, it rises from \bneNsmePrePh{} to \bneNsmePostPh{}.
These are large changes relative to the pre-policy protected pools. Non-SME
participation moves in the opposite direction, falling from \bneNnsPreNp{} to
\nonSmePostNp{} in non-pharma and from \bneNnsPrePh{} to \nonSmePostPh{} in
pharma.

That participation response matters for interpretation. If one held the SME
pool fixed at its pre-policy size, the set-aside would look more expensive than
the realized post-policy environment. The government gets some competition back
through induced SME participation. The key question is whether that offset is
large enough to replace the price discipline supplied by the excluded non-SMEs.

\subsection{Exclusion dominates}

Table~\ref{tab:price_decomp_v8} reports the three simulated price statistics.
The open pre-policy auction is $S_1$. The SME-only auction holding the
pre-policy SME pool fixed is $S_2$. The SME-only auction with the observed
post-policy SME pool is $S_3$. Prices are normalized by the buyer's reference
price and correspond to the simulated second-order statistic, $E[c_{(2)}]$.

The first move, $S_1\rightarrow S_2$, isolates bidder exclusion. In
non-pharmaceutical markets, removing non-SMEs while holding the pre-policy SME
pool fixed raises the simulated price from \bneMeanSoneNp{} to
\bneMeanStwoNp{}, an increase of \bneEffectIntensNp{} of the reference price.
In pharmaceutical markets, the same move raises the simulated price from
\bneMeanSonePh{} to \bneMeanStwoPh{}, an increase of \bneEffectIntensPh{}.
This is the price-forming cost of replacing the full bidder pool with the
protected pool.

The second move, $S_2\rightarrow S_3$, allows the post-policy SME pool to
enter. This term is negative in both classes. In non-pharma, SME
entry/composition reduces the simulated price by \bneEffectEntryNp{} of the
reference price. In pharma, it reduces the simulated price by
\bneEffectEntryPh{}. The realized set-aside effect, $S_3-S_1$, remains
positive: \bneEffectTotalNp{} in non-pharma and \bneEffectTotalPh{} in pharma.
Entry attenuates the price shock; it does not eliminate it.

\begin{table}[!htbp]
\centering
\begin{threeparttable}
\caption{Price decomposition: exclusion versus entry/composition}
\label{tab:price_decomp_v8}
\footnotesize
\setlength{\tabcolsep}{4pt}
\begin{tabular}{lccccccc}
\toprule
 & \multicolumn{3}{c}{Simulated price} & \multicolumn{4}{c}{Decomposition} \\
\cmidrule(lr){2-4}\cmidrule(lr){5-8}
Class & $S_1$ & $S_2$ & $S_3$ &
$S_2-S_1$ & $S_3-S_2$ & $S_3-S_1$ & Excl./net \\
\midrule
Non-pharma
& \bneMeanSoneNp{} & \bneMeanStwoNp{} & \bneMeanSthreeNp{}
& +\bneEffectIntensNp{} & \bneEffectEntryNp{}
& +\bneEffectTotalNp{} & \bneIntensRelNetNp{}\% \\
Pharma
& \bneMeanSonePh{} & \bneMeanStwoPh{} & \bneMeanSthreePh{}
& +\bneEffectIntensPh{} & \bneEffectEntryPh{}
& +\bneEffectTotalPh{} & \bneIntensRelNetPh{}\% \\
\bottomrule
\end{tabular}
\begin{tablenotes}\footnotesize
\item $S_1$ is the open auction with the pre-policy bidder pool. $S_2$ is the
SME-only auction holding the pre-policy SME pool fixed. $S_3$ is the SME-only
auction using the observed post-policy SME pool. Prices are simulated
$E[c_{(2)}]$ values normalized by the buyer reference price. ``Excl./net''
reports $(S_2-S_1)/(S_3-S_1)$. The entry/composition term is negative because
the post-policy protected pool partially offsets the exclusion shock.
\end{tablenotes}
\end{threeparttable}
\end{table}

The relative-to-net column is intentionally reported because it shows why the
absolute share can understate the size of the exclusion margin. In non-pharma,
the exclusion component is \bneIntensRelNetNp{} percent of the net effect; in
pharma it is \bneIntensRelNetPh{} percent. The entry/composition term offsets
\bneEntryRelNetNp{} percent and \bneEntryRelNetPh{} percent of the net effect,
respectively. If instead one normalizes by the sum of absolute component
magnitudes, exclusion accounts for \bneShareIntensNp{} percent in non-pharma
and \bneShareIntensPh{} percent in pharma. Both normalizations point in the
same direction: the dominant force is the loss of the non-SME price-forming
pool, not a lack of SME entry.

Figure~\ref{fig:decomposition_v8} makes the same point visually. The jump from
$S_1$ to $S_2$ is the cost of exclusion. The decline from $S_2$ to $S_3$ is the
offset from the realized protected pool. The remaining gap between $S_1$ and
$S_3$ is the realized price cost of the set-aside under observed entry.

\begin{figure}[!htbp]
\centering
\includegraphics[width=0.95\textwidth]{fig_v3_decomposition.pdf}
\caption[Counterfactual price decomposition.]{\textit{Counterfactual price
decomposition.} Bars report the simulated second-order statistic normalized by
the buyer reference price. $S_1$ is the open pre-policy auction, $S_2$ removes
non-SMEs while holding the pre-policy SME pool fixed, and $S_3$ allows the
observed post-policy SME pool to enter. The movement from $S_1$ to $S_2$ is the
lost competitive-discipline component; the movement from $S_2$ to $S_3$ is the
entry/composition offset.}
\label{fig:decomposition_v8}
\end{figure}

\subsection{Entry and composition are offsets, not the source}

The decomposition separates two facts that are often conflated. The set-aside
does induce protected entry. That entry is valuable in price terms because it
pulls the simulated price down from $S_2$ to $S_3$. But the policy first creates
a larger price increase by removing non-SMEs from the auction. In the baseline
simulation, holding the protected pool at its pre-policy size would make the
latent price shock \latentShockRatioNp{} percent larger in non-pharma and
\latentShockRatioPh{} percent larger in pharma. The realized post-policy SME
pool absorbs part of that shock, but the net effect remains positive in both
classes.

The entry/composition term should not be overinterpreted as pure entry. It also
contains whatever empirical difference exists between the active SME bidders
observed before and after the cutoff. That is why the paper treats
$F_c^{\mathrm{SME,Post}}$ explicitly as an observed post-policy active-bidder
distribution rather than as the output of a formal selection model. Two
diagnostics keep the conclusion focused. First, the decomposition remains
exclusion-dominant under the Turnbull cost estimator, with exclusion shares of
\turnbullShareNp{} percent in non-pharma and \turnbullSharePh{} percent in
pharma on the absolute-share normalization. Second, under the strict-invariance
benchmark that forces the post-policy SME distribution to equal the pre-policy
one, the exclusion component remains the larger component
(\strictShareNp{} percent in non-pharma and \strictSharePh{} percent in
pharma). What changes under strict invariance is the policy ranking in pharma,
not the basic price-formation fact that bidder exclusion is the larger force.

Firm turnover helps explain why the pharma ranking is more fragile. In
pharmaceuticals, \turnoverPhPctNewFirms{} percent of post-policy SME firms are
new relative to the pre-period pool and they account for
\turnoverPhPctNewBids{} percent of post-policy SME bids. In non-pharma, the
corresponding bid share is \turnoverNpPctNewBids{} percent. The composition
shift is therefore larger precisely in the class where the welfare ranking is
model-sensitive. This is a scope condition, not a competing headline. The
robust result is that in the thicker non-pharmaceutical market, the price cost
is driven by removing non-SMEs and survives alternative treatments of the
protected pool.

\subsection{Interpretation}

The mechanism is not that protected firms strategically bid less aggressively
after the set-aside. In the maintained Preg\~ao/IPV model, bidders exit at
cost. The mechanism is that the set-aside changes which cost order statistic
forms the price. In the open auction, non-SMEs remain in the pool and discipline
the second-lowest cost. In the SME-only auction, that discipline is removed.
Additional SME entry pushes back against the loss, but the post-policy
protected pool does not recreate the full-pool benchmark.

This is the price-formation result the welfare section builds on. A policy
that supports SMEs by excluding non-SMEs buys redistribution by weakening the
auction's price-forming pool. A policy that supports SMEs while leaving
non-SMEs in the auction works through a different channel: it can tilt
allocation toward SMEs while preserving the bidders that discipline the price.
